\section{Related Work}
\label{sec:related}
\textbf{VoE benchmarks.} IntPhys pairs possible and impossible videos rendered in Unreal Engine, matched pixel for pixel, covering object permanence, shape constancy and continuity \cite{riochet2022}. InfLevel provides VoE experiments for continuity, solidity and gravity, on which the video architectures it tested were at or near chance \cite{weihs2022}. IntPhys~2 uses more complex scenes; V-JEPA and V-JEPA~2 are near chance on it (the best predictive model, V-JEPA~2, reaches 57.5\%), and a qualitative debug scene shows no extra surprise at the point where an occluded ball should reappear \cite{bordes2025}. Physion and Physion++ take another route: models predict object contact (and, in Physion++, infer latent properties such as mass and friction) and are compared with humans \cite{bear2021,tung2023}. We evaluate on IntPhys because its matched pairs make the pairwise verdict well defined.

\textbf{Latent prediction and intuitive physics.} V-JEPA learns video representations by predicting masked regions in representation space, with no pixel reconstruction \cite{bardes2024}; V-JEPA~2 scales the recipe to more video and adds a robot-planning stage \cite{assran2025}. Garrido et al.\ used V-JEPA's prediction error as a surprise signal and reported 98\% pairwise accuracy on IntPhys dev for their ViT-H model (context length tuned per property), with pixel-prediction and multimodal language-model baselines near chance \cite{garrido2025}. VideoMAE~V2, a masked autoencoder that reconstructs pixels, was planned as a pixel-reconstruction contrast model \cite{wang2023}; we did not run it. Joseph et al.\ probe video encoders layer by layer for physical information and find it concentrated at intermediate layers \cite{joseph2026}. These works report that the signal scores well on one benchmark. We ask what it responds to, for a public checkpoint.

\textbf{Doubts about what accuracy shows.} Su et al.\ report near-chance IntPhys~2 accuracy (0.50--0.54) for V-JEPA models trained on three different video sources, and note qualitatively that frame-wise surprise was often similar for possible and impossible videos \cite{su2025}. Li et al.\ find that on Physion, frozen V-JEPA~2 features with probes reach accuracy close to humans, while its disagreement with humans on individual items is far higher than between humans \cite{li2026}. That is a contact-prediction task, not a VoE one. Both point to the same gap: accuracy alone does not say what the model uses. We add one control for the VoE setting, abrupt edits to possible clips whose pixel magnitude is matched to the real violations.

\textbf{Counterfactual and likelihood-based evaluation of video models.} CRONOS probes whether video generators stay physically consistent under controlled visual interventions such as viewpoint, scene and appearance, and measures generation quality, not model surprise \cite{begiristain2026}. Physics-IQ evaluates the physical understanding of generative video models \cite{motamed2025}. LikePhys compares the denoising loss of matched valid and invalid videos for diffusion models \cite{likephys2025}. YoCausal treats denoising loss as surprise and uses time reversal of real videos as a counterfactual for diffusion models \cite{yocausal2026}. We use time reversal as one positive control; those results do not tell us whether it works for latent predictors, so we test it (Section~\ref{sec:method}). Our edits target the surprise of frozen predictive encoders.

\textbf{Impossible versus improbable.} Michaelov et al.\ found that language models do not reliably assign higher probability to possible-but-improbable sentences than to impossible ones \cite{michaelov2025}. Our specificity test asks the analogous question for video surprise.
